% =====================================================================
%  NO4Cheetah -- First draft for supervisor feedback (v1, 29.09.2026)
%  Snippets to paste into the IEEE template (E-10, TUHH).
%
%  Needed in the preamble (if not already there):
%    \usepackage{amsmath,amssymb}
%    \usepackage{xcolor}
%    \newcommand{\ph}[1]{\textcolor{red}{[#1]}}   % placeholder marker
%
%  Placeholder convention: \ph{WERT}, \ph{BILD}, \ph{FORMEL}, \ph{ZITAT},
%  \ph{TABELLE}, \ph{ABSCHNITT}. German %-comments are notes to the author.
%  Bib keys used are listed in refs_draft.bib (check all entries).
% =====================================================================


% =====================================================================
\begin{abstract}
% 150-250 words, no symbols/maths/citations. Currently ~230 words incl.
% placeholders. Results sentences are placeholders -- fill after Sec. 4.
Differentiable beam dynamics codes such as Cheetah enable gradient-based
tuning and system identification of particle accelerators. When
three-dimensional space charge is included, the grid-based field
calculation accounts for a large share of the computation time and of the
memory required for reverse-mode differentiation. In each space charge
kick, the
charge density of the beam is convolved with an integrated Green function
by fast Fourier transforms on a grid that is doubled in every dimension to
represent open boundaries. This work replaces this field solve with a
Fourier neural operator, a neural network that learns the mapping from
charge density to electrostatic potential as an operator between function
spaces and shares the Fourier-space convolution structure of the classical
solver. The operator is trained on data generated with the existing solver
and embedded in the kick without changing its interface, so that it remains
usable in arbitrary beamlines. Computation time and peak memory of the
forward and backward pass are compared with the Green function solver at
equal and finer grid resolution on CPU and GPU. \ph{KERNBEFUND: e.g. The
surrogate reduces the time of the field solve by about \ph{WERT} and the
recorded memory by \ph{WERT}, at a relative error of the potential of
\ph{WERT} percent.} \ph{GRENZE: e.g. Charge deposition and force
interpolation, which scale with the number of particles, remain unchanged
and limit the overall gain for large particle numbers.}
\end{abstract}


% =====================================================================
\section{Introduction}
\label{sec:intro}

% --- 2.a Aim of this work / research question -------------------------
% Hinweis: Viele Leser brauchen den Kontext aus 2.b vor der Forschungsfrage.
% Alternative: diesen Absatz ans Ende der Einleitung stellen.
Cheetah~\cite{kaiser2024cheetah} is a differentiable beam dynamics code
whose three-dimensional space charge kick~\cite{dhamrait2026memory} is
costly in computation time and in the memory needed for reverse-mode
differentiation (Sec.~\ref{sec:intro-cost}). For $10^4$ macroparticles on
the default grid of $32^3$ cells, most of this cost arises in the
grid-based field solve, which maps the charge density of the beam to its
electrostatic potential. This work investigates whether a Fourier neural
operator (FNO)~\cite{li2021fno} can replace this field solve such that the
computation time and the memory of the forward and backward pass decrease,
while the relative error of the potential stays below \ph{WERT}\,\% and the
beam moments after \ph{WERT} kicks deviate by less than \ph{WERT}\,\% from
the reference solver.
% Optional, falls Ziel der Arbeit (mit Betreuern klaeren):
% "... and the gradients of beam properties with respect to the input
%  parameters agree with those of the reference solver (cosine similarity
%  above \ph{WERT})."


% --- 2.b.i Cheetah: related work --------------------------------------
\subsection{Cheetah and space charge}
\label{sec:intro-cheetah}

Cheetah~\cite{kaiser2024cheetah} is a beam dynamics code written in
PyTorch. Lattice elements are modular objects that map an incoming to an
outgoing beam, so that arbitrary beamlines are assembled by composing
elements. Tensorised computation shortens simulations by orders of
magnitude compared with established codes; in one reinforcement-learning
application, training time dropped from over 12~days with
Ocelot~\cite{agapov2014ocelot} to just over one hour with
Cheetah~\cite{kaiser2024cheetah}. Since all operations are PyTorch
operations, gradients of beam properties with respect to lattice and beam
parameters follow from automatic differentiation, which enables
gradient-based tuning and system identification~\cite{kaiser2024cheetah}.
Dhamrait et al.~\cite{dhamrait2026memory} added a three-dimensional space
charge kick based on the particle-in-cell (PIC) method; this kick is the
subject of this work.

% --- 2.b.ii Role of space charge in beamline simulations -------------
Space charge denotes the Coulomb interaction among the particles of a
beam. Its self-field defocuses the beam, with the magnetic contribution
partially cancelling the electric one by a factor $1/\gamma^2$, where
$\gamma$ is the Lorentz factor~\cite{schindl2006spacecharge}. It is
therefore most pronounced for high bunch charge at low energy and drives,
e.g., emittance growth and intensity-dependent beam
loss~\cite{dhamrait2026memory}. Because the force depends on the
distribution of the whole beam, space charge cannot be written as a
transfer map of the element. It is applied instead as a sequence of
kicks: between two kicks the beam is transported by the element map, and
each kick changes the momentum of every macroparticle by the self-field
force integrated over the time $\Delta t$ in which the beam travels the
step $\Delta s$~\cite{dhamrait2026memory}. A direct summation over all
particle pairs scales as $\mathcal{O}(N^2)$ for $N$ macroparticles. PIC
methods instead deposit the charge on a grid, solve for the field on the
grid, and interpolate the force back to the
particles~\cite{birdsall2004plasma,hockney1988computer}. A kick thus
consists of a particle part, whose cost grows with $N$, and a grid part,
whose cost grows with the number of cells $N_\mathrm{cells}$
(Sec.~\ref{sec:method-placement}).

% --- 2.b.iii Computation time and memory ------------------------------
% Zahlen unten: VORLAEUFIG aus P2-Profiling (21.09.2026, cheetah 0.8.4,
% CPU 4 Kerne, float32). Vor Versand durch finale Messung ersetzen oder
% im Text als vorlaeufig kennzeichnen. Messaufbau: Sec. method-measurement.
\label{sec:intro-cost}
Table~\ref{tab:profile} lists the computation time of one kick in
Cheetah, broken down by step. For $10^4$ macroparticles on $32^3$ cells,
the forward pass takes $46.0\,\mathrm{ms}$ and the backward pass
$41.5\,\mathrm{ms}$. Of the forward time, $28.7\,\mathrm{ms}$ ($62\,\%$)
are spent evaluating the integrated Green function (IGF) and
$3.2\,\mathrm{ms}$ ($7\,\%$) on the FFT-based convolution; on $64^3$ cells
the IGF share rises to $78\,\%$. The IGF depends only on the cell size and
$\gamma$, not on $N$. Charge deposition and force interpolation, in
contrast, grow with $N$ and account for $60\,\%$ of the forward time at
$10^5$ macroparticles on $32^3$ cells. For comparison, a linear element
tracks the same beam in \ph{WERT}\,ms.
% ^ Messung nachtragen (z.B. Drift oder Quadrupol, gleiche N) -- macht
%   "teuer" erst messbar.
Memory is the second constraint. For the backward pass, PyTorch stores
intermediate results of every kick until the gradient is computed.
Dhamrait et al.~\cite{dhamrait2026memory} measured that this memory grows
linearly with $N$ and $N_\mathrm{cells}$ and proportionally to the number of
kicks, with $1.4\,\mathrm{kB}$ per macroparticle, dominated by deposition
and interpolation, and $1.2\,\mathrm{kB}$ per cell, dominated by the IGF
evaluation. For $10^4$ macroparticles on $32^3$ cells this amounts to
about $53\,\mathrm{MB}$ per kick, i.e.\ \ph{WERT}\,GB for a beamline with
\ph{WERT} kicks. Because time and memory of the grid part grow with
$N_\mathrm{cells}$, the grid resolution trades cost against the accuracy
of the field \ph{ZITAT -- Beleg fuer: Genauigkeit vs. Gitteraufloesung;
oder eigene Aufloesungsstudie, Abschnitt \ph{ABSCHNITT}}.

\begin{table}[t]
  \centering
  \caption{Forward and backward computation time of one space charge kick
  in Cheetah~\cite{kaiser2024cheetah,dhamrait2026memory}, in ms, broken
  down by step (forward pass). CPU \ph{WERT: Modell, 4 Kerne}, float32,
  Cheetah \ph{WERT}, PyTorch \ph{WERT}. IGF: evaluation of the integrated
  Green function on the doubled grid; FFT: convolution; Dep.+Int.: charge
  deposition and force interpolation. \ph{WERT: Median ueber n
  Wiederholungen, Streuung}.}
  \label{tab:profile}
  \small
  \setlength{\tabcolsep}{4pt}
  \begin{tabular}{rrrrrrr}
    \hline
    Grid & $N$ & IGF & FFT & Dep.+Int. & Fwd & Bwd \\
    \hline
    $32^3$ & $10^4$ &  28.7 &  3.2 &  8.6 &  46.0 &  41.5 \\
    $32^3$ & $10^5$ &  20.7 &  1.7 & 57.7 &  95.8 & 120.3 \\
    $64^3$ & $10^4$ & 136.8 & 22.7 &  8.5 & 174.4 & 134.4 \\
    $64^3$ & $10^5$ & 116.0 & 22.7 & 85.0 & 240.8 & 216.2 \\
    \hline
  \end{tabular}
\end{table}
% Optional statt/zusaetzlich zur Tabelle: \ph{BILD} gestapelte Balken je
% Schritt, vorwaerts/rueckwaerts nebeneinander (vollstaendige Plots -> Anhang).


% --- 2.c.i Fourier neural operators: related work ---------------------
\subsection{Fourier neural operators}
\label{sec:intro-fno}

Neural operators learn mappings between function spaces instead of
between finite-dimensional
vectors~\cite{kovachki2023neural,li2021fno,azizzadenesheli2024neural}. In
the present case, the input function is the charge density $\rho$ and the
output function the potential $\phi$; the solution operator
$\mathcal{G}\colon\rho\mapsto\phi$ is linear and, for open boundaries,
given by the convolution with the Green function
(Eq.~\eqref{eq:conv}). A conventional network, e.g.\ a convolutional
network, approximates this map for one fixed grid, and its error grows
when the resolution changes~\cite{li2021fno}. A neural operator
parametrises $\mathcal{G}$ itself; its parameters do not depend on the
discretisation, so one trained model can be evaluated on different grids,
provided the grid resolves the relevant
features~\cite{li2021fno,kovachki2023neural}. The FNO~\cite{li2021fno}
lifts the input pointwise to $d_v$ channels, applies $L$ Fourier layers
\begin{equation}
  v_{t+1}(\mathbf{x}) = \sigma\!\left( W v_t(\mathbf{x})
  + \mathcal{F}^{-1}\!\big( R \cdot \mathcal{F}(v_t) \big)(\mathbf{x})
  \right),
  \label{eq:fourier-layer}
\end{equation}
and projects the result pointwise to the output. Here $\mathcal{F}$ is the
Fourier transform, $R$ a learned complex weight tensor acting on the lowest
$k_\mathrm{max}$ modes per dimension, $W$ a learned pointwise linear map,
and $\sigma$ a nonlinear activation. With FFTs, one layer costs
$\mathcal{O}(n\log n)$ for $n$ grid points~\cite{li2021fno}. Li et
al.~\cite{li2021fno} report consistent errors across resolutions and
evaluation on finer grids than used in training (zero-shot
super-resolution).

% --- 2.c.ii Why an FNO for this work ----------------------------------
The FNO fits the space charge problem because it shares the structure of
the classical solver. Cheetah computes the potential as
$\phi = \mathcal{F}^{-1}\big(\mathcal{F}(\tilde G)\cdot
\mathcal{F}(\tilde\rho)\big)$ on a doubled grid (Eq.~\eqref{eq:hockney}),
while a Fourier layer computes $\mathcal{F}^{-1}\big(R\cdot
\mathcal{F}(v)\big)$. A single Fourier layer on the doubled grid, without
nonlinearity, local term and mode truncation and with
$R=\mathcal{F}(\tilde G)$, thus reproduces the solver exactly; the FNO
contains the classical solver as a special case. This has three
consequences for Cheetah. First, the learned weights replace the
evaluation of the IGF, which dominates time and memory of the kick for
moderate $N$ (Sec.~\ref{sec:intro-cost}). Second, discretisation
invariance allows a refinement of the \emph{grid} without retraining; it
does not affect the particle part of the kick, whose cost grows with $N$.
Third, Li et al.\ attribute to the local term $W$ the ability to represent
non-periodic boundary conditions~\cite{li2021fno}, shown for problems with
bounded domains. Whether this extends to the open boundaries of the space
charge problem, so that the grid doubling of the classical solver can be
omitted, is tested in this work (Sec.~\ref{sec:method-fno}).
This approach differs from the space charge surrogate of Kaiser et
al.~\cite{kaiser2024cheetah}: their network predicts the space-charge
induced change of the second beam moments in a quadrupole, which takes
$370\,\mu\mathrm{s}$ compared with $1.36\,\mathrm{s}$ for tracking with
space charge in Ocelot, but acts on moment-based beams and is tied
to one element. The FNO replaces the field solve inside the kick, keeps
the element interface, and therefore acts on particle beams in any
beamline.


% =====================================================================
\section{Methodology}
\label{sec:method}

\subsection{FNO surrogate in Cheetah}

% --- 3.a.iii Placement in Cheetah -------------------------------------
% Reihenfolge-Vorschlag: iii (Einbettung) vor i (Aufbau), weil der Aufbau
% erst mit bekannten Ein-/Ausgaengen nachvollziehbar ist.
\label{sec:method-placement}
A space charge kick over the step $\Delta s$ consists of the following
steps~\cite{dhamrait2026memory}:
\begin{enumerate}
  \item The macroparticles are transformed from coordinates at fixed $s$
  to coordinates at fixed time $t$.
  \item Their charge is deposited on a Cartesian grid of $32^3$ cells
  spanning $\pm3\sigma$ of the beam in each dimension by trilinear
  (cloud-in-cell) weighting~\cite{birdsall2004plasma}, which yields the
  charge density $\rho_{ijk}$.
  \item The potential follows from the modified Poisson
  equation~\cite{vay2008simulation}
  \begin{equation}
    \partial_x^2\phi + \partial_y^2\phi + \frac{1}{\gamma^2}\partial_z^2\phi
    = -\frac{\rho}{\epsilon_0},
    \label{eq:poisson}
  \end{equation}
  solved in the laboratory frame by convolution with the
  IGF~\cite{qiang2006three,dhamrait2026memory}
  \begin{equation}
    \phi_{ijk} = \sum_{i'j'k'} \rho_{i'j'k'}\,
    G(\mathbf{x}_{ijk}-\mathbf{x}_{i'j'k'}).
    \label{eq:conv}
  \end{equation}
  $G$ is the Green function of Eq.~\eqref{eq:poisson} integrated over one
  cell. It is evaluated in closed form from an antiderivative $g$ at the
  eight corners of the cell~\cite{qiang2024implementation},
  \begin{equation}
    G(\mathbf{x}) = \sum_{a,b,c\in\{\pm1\}} abc\;
    g\!\left(x+a\tfrac{\Delta x}{2},\, y+b\tfrac{\Delta y}{2},\,
    \gamma z+c\tfrac{\gamma\Delta z}{2}\right).
    \label{eq:igf}
  \end{equation}
  Since Eq.~\eqref{eq:conv} costs $\mathcal{O}(N_\mathrm{cells}^2)$ when
  evaluated directly, Cheetah uses the Hockney
  method~\cite{hockney1988computer}: $\rho$ is zero-padded and $G$ is
  evaluated on a grid doubled in every dimension ($\tilde\rho$, $\tilde G$,
  $64^3$ points), so that the cyclic convolution computed by the FFT equals
  the linear convolution of Eq.~\eqref{eq:conv}, i.e.\ open boundaries:
  \begin{equation}
    \phi = \mathcal{F}^{-1}\!\big(\mathcal{F}(\tilde G)\cdot
    \mathcal{F}(\tilde\rho)\big)\Big|_{32^3}.
    \label{eq:hockney}
  \end{equation}
  \item The force $\mathbf{F} = -q\nabla\phi/\gamma^2$ is computed on the
  grid by central differences~\cite{dhamrait2026memory}.
  \item $\mathbf{F}$ is interpolated trilinearly to the macroparticles.
  \item Each momentum is updated, $\mathbf{p}\leftarrow\mathbf{p}+
  \mathbf{F}\Delta t$ with $\Delta t = \Delta s/(\beta c)$, and the
  coordinates are transformed back to fixed $s$.
\end{enumerate}
The surrogate replaces step~3, i.e.\ the evaluation of $\tilde G$ and the
convolution of Eq.~\eqref{eq:hockney}, and maps $\rho$ on $32^3$ cells to
$\phi$ on the same grid. Steps~1, 2 and 4--6 remain unchanged, and so does
the interface of the kick element; the surrogate kick can therefore be
placed in any Cheetah beamline. The replaced routine is given in
Listing~\ref{lst:solve} in the appendix.

% --- 3.a.i Structure of the FNO ----------------------------------------
\label{sec:method-fno}
The input of the FNO is the charge density of the whole beam on the grid,
$\rho\in\mathbb{R}^{B\times32\times32\times32}$ for a batch of $B$ beams.
Because $\tilde G$ depends on the cell size $(\Delta x,\Delta y,
\gamma\Delta z)$ (Eq.~\eqref{eq:igf}), the operator $\rho\mapsto\phi$
depends on these parameters as well, and the FNO must receive them. We
remove the dependence on scale by normalising
\begin{equation}
  \hat\rho = \frac{\rho L^3}{Q}, \qquad
  \hat\phi = \frac{4\pi\epsilon_0 L}{Q}\,\phi,
  \label{eq:nondim}
\end{equation}
with the total charge $Q$ on the grid and the mean edge length $L$ of the
grid box \ph{FORMEL: Definition von L, inkl. $\gamma$-Streckung}, and pass
the remaining aspect ratios of the box as constant input channels.
% ^ nur stehen lassen, wenn so umgesetzt. Kalibrierlauf: L an Box statt
%   Zelle gekoppelt -> 3 % statt 142 % Fehler beim Wechsel 16^3 -> 32^3
%   (gehoert dann in Ergebnisse, nicht hierher).
The FNO lifts the input to $d_v=\ph{WERT}$ channels, applies $L_F=\ph{WERT}$
Fourier layers (Eq.~\eqref{eq:fourier-layer}) with
$k_\mathrm{max}=\ph{WERT}$ modes per dimension and \ph{WERT: Aktivierung},
and projects to one output channel, $\hat\phi\in
\mathbb{R}^{B\times32\times32\times32}$; the model has \ph{WERT} parameters.
The number of layers is a hyperparameter, set to \ph{WERT} following
\ph{ZITAT -- Details zu: Hyperparameterwahl, z.\,B. Li et al. (vier
Schichten)} or a hyperparameter search (\ph{ABSCHNITT}).
Two boundary treatments are compared: (a) the input zero-padded to $64^3$
as in the Hockney method, and (b) the unpadded $32^3$ grid, relying on the
local term $W$ for the open boundaries~\cite{li2021fno}. Variant~(b) uses
one eighth of the grid points of variant~(a).
% Randbedingungen, die unabhaengig vom Surrogat bleiben (ggf. 1 Satz):
% zentrale Differenzen setzen den Gradienten in der aeussersten Zellschicht
% auf null, Teilchen ausserhalb +-3 sigma erhalten keine Kraft.

% --- 3.a.ii Training and evaluation -----------------------------------
The training data are pairs $(\hat\rho,\hat\phi)$ generated with the IGF
solver of Cheetah for \ph{WERT} beams drawn from \ph{WERT: Verteilungen,
z.\,B. gaussisch und nicht-gaussisch} with $\gamma\in\ph{WERT}$ and
\ph{WERT: weitere Parameterbereiche}, split \ph{WERT} into training,
validation and test sets. The loss is the relative $L_2$ error of
$\hat\phi$ plus \ph{WERT} times the relative $L_2$ error of its
central-difference gradient, because the kick depends on $\nabla\phi$ and
differentiation amplifies high-frequency errors. We train with Adam,
learning rate \ph{WERT}, batch size \ph{WERT}, for \ph{WERT} epochs on
\ph{WERT: Hardware}.
The model is evaluated on three levels against the IGF solver at the same
and at finer resolution: the relative $L_2$ error of $\phi$ and
$\nabla\phi$, reported as median, 95th percentile and maximum over the
test set; the deviation of beam sizes and emittances after \ph{WERT}
kicks; and \ph{falls Ziel: the cosine similarity between the gradients of
a beam objective with respect to the input parameters computed through
surrogate and solver}. Resolution transfer is tested by evaluating a
model trained on $32^3$ cells on $64^3$ cells without retraining.
% Optional (empfohlen, sonst ist "Operator statt CNN" nur behauptet):
% Ablation gegen ein CNN/U-Net mit gleicher Parameterzahl auf festem Gitter.

% --- 3.b Measurement of computation time and memory -------------------
\subsection{Measurement of computation time and memory}
\label{sec:method-measurement}
Three field solvers are compared inside the otherwise unchanged kick:
the IGF solver as implemented, the IGF solver with $\tilde G$ computed once
and reused, since $\tilde G$ depends only on cell size and $\gamma$
(Eq.~\eqref{eq:igf}), and the FNO. For each solver, we measure time and
peak memory of one kick, forward and backward pass separately, and break
the time down by the steps of Sec.~\ref{sec:method-placement}. Wall time
is the median of \ph{WERT} repetitions after one warm-up run, measured
with the PyTorch benchmark utility and with device synchronisation on the
GPU. Peak memory is read from the allocator on the GPU and as the maximum
resident set size of a separate process per configuration on the CPU. We
vary the grid ($32^3$, $64^3$) and $N$ ($10^4$, $10^5$, $10^6$) on a
\ph{WERT: CPU} and a \ph{WERT: GPU}, with PyTorch \ph{WERT}, Cheetah
\ph{WERT} and float32. Speed-ups are reported at equal accuracy of the
potential, not only at equal resolution. Code and environment
(\texttt{requirements.txt}) are available at \ph{ABSCHNITT: TUHH-GitLab}.


% =====================================================================
\section{Results}
\label{sec:results}
% 4.a Zeit und Speicher FNO vs. Status quo. Muster je Abbildung:
%   (1) was dargestellt ist, (2) Beobachtung mit Zahlen, (3) welche
%   Aussage das stuetzt. Vorlage:
Fig.~\ref{fig:timing} shows forward and backward time per kick over the
number of cells for the IGF solver, the IGF solver with reused $\tilde G$,
and the FNO. \ph{WERT: Beobachtung, z.\,B. ab $n^3$ Zellen ist das FNO
um Faktor x schneller}. \ph{Aussage, die das stuetzt}.
% 4.b Was das FNO leistet: Fehler Feld/Strahl(/Gradient), Aufloesungstransfer.
% 4.c Was fehlt: z.B. Randfehler, Gradiententreue, Teilchenteil unveraendert.


% =====================================================================
\section{Conclusion}
\label{sec:conclusion}
% 5.a Kernbefunde, 5.b Antwort auf die Forschungsfrage (gut genug? was
% gewinnt Cheetah?) -- nach den Ergebnissen schreiben.

% --- 5.c Outlook -------------------------------------------------------
The surrogate replaces only the grid part of the kick. Charge deposition
and force interpolation scale with $N$; they dominate the memory per
macroparticle in reverse mode~\cite{dhamrait2026memory} and, in our
profile, the computation time at $10^5$ macroparticles
(Table~\ref{tab:profile}). Reducing this part requires a model that acts on
the particle cloud directly, e.g.\ a permutation-invariant network such as
PointNet~\cite{qi2017pointnet}, which Kaiser et al.\ proposed for
extending Cheetah's surrogates to particle beams~\cite{kaiser2024cheetah}.
Such a model maps a point cloud instead of a function on a grid and thus
gives up discretisation invariance. Independently of the surrogate,
gradient checkpointing reduces the memory of reverse mode at the cost of
recomputation~\cite{dhamrait2026memory}.


% =====================================================================
\appendix
\section{Code}
% Listing der ersetzten Routine (Original + Surrogat-Variante)
\begin{figure}[h]
  \centering
  \ph{Listing: Poisson-Loesung mit FNO statt IGF-Faltung}
  \caption{Replaced field solve of the space charge kick.}
  \label{lst:solve}
\end{figure}
% -> mit listings/minted als lstlisting-Umgebung ersetzen.

\section{Computation time and memory measurements}
% Alle Plots aus Sec. method-measurement; im Haupttext nur die
% Kernabbildung(en).
\begin{figure}[h]
  \centering
  \ph{BILD: Zeit je Kick ueber Zellzahl, alle Konfigurationen}
  \caption{\ph{Bildunterschrift}}
  \label{fig:timing}
\end{figure}
